% chapter13.tex -- Complexity Analysis
% Based on the author's UGThesis.docx.

\chapter{Complexity Analysis}
\label{ch:complexity}

In this chapter we analyse the time and space complexity of our algorithm. 
We show that the total running time is \(O(T)\), where \(T\) is the number of triangulations generated, which yields \(O(1)\) amortised time per triangulation. 
We also prove that the space complexity is \(O(n)\), where \(n\) is the number of vertices of the input graph.

\section{Cost Model}
\label{sec:cost_model}

We count elementary operations that can be performed in \(O(1)\) time:
\begin{itemize}
    \item Insertion, deletion, and membership test in the global hash set \(S\).
    \item Insertion and removal of a chord from the linked lists \(L\), \(GS\), and \(VGS\).
    \item Updating the opposite pair of a generating chord.
    \item Flipping a chord (removing one chord, adding another).
\end{itemize}
All these operations are assumed to take \(O(1)\) expected time (for hash set operations) or worst-case \(O(1)\) (for list operations).

Let \(B\) denote the total number of \textbf{building operations} (chord insertions when constructing root triangulations) and \(F\) the total number of \textbf{flipping operations} (edge flips that generate child triangulations) across the entire execution. 
Since each operation costs \(O(1)\), the total running time is \(O(B + F)\). 
We will prove that \(B + F = O(T)\), where \(T\) is the total number of triangulations output. 
Hence the amortised time per triangulation is \(O(1)\).

\section{Minimum Triangulations Guarantee}
\label{sec:min_triangulations}

Before analysing the total work, we recall a crucial lower bound established in Chapter~\ref{ch:lower_bound}.

\begin{lemma}[Minimum triangulations per face]
    \label{lem:min_triangulations_final}
    For any face \(F_i\) with \(n_i\) vertices, the number of valid triangulations \(d_i\) that are compatible with the global chord set \(S\) (from previous faces) satisfies
    \[
        d_i \ge n_i - 3.
    \]
    Consequently, \(d_i = \Omega(n_i)\).
\end{lemma}

This lemma is proved in Theorem~\ref{thm:lower_bound} by considering two safe root triangulations and the path between them in the genealogical tree. 
The bound is tight in the worst case (e.g., a quadrilateral with one diagonal already blocked yields exactly one triangulation, and \(n_i-3 = 1\); a pentagon can have as few as two triangulations, etc.). 
The important consequence is that \(n_i = O(d_i)\), i.e., the number of vertices of a face is at most a constant times the number of triangulations generated for that face.

\section{Building Operations Analysis}
\label{sec:building_ops}

For a face \(F_i\) with \(n_i\) vertices, building its safe root triangulation involves:
\begin{itemize}
    \item Finding a safe root using the two-pointer algorithm: \(O(n_i)\) time.
    \item Adding \(n_i-3\) chords to form the fan triangulation: \(O(n_i)\) time.
    \item Initialising the valid generating set (VGS) by checking the opposite pairs of the \(n_i-3\) chords against \(S\): \(O(n_i)\) time.
\end{itemize}
Thus the total building cost for a single invocation of the root of face \(F_i\) is \(O(n_i)\).

However, in the overall algorithm, the root triangulation of a face may be built many times: once for each distinct combination of triangulations of the previous faces. 
Let \(T_{i-1} = \prod_{j=1}^{i-1} d_j\) be the number of ways to triangulate faces \(F_1, \dots, F_{i-1}\). 
Then the root triangulation of \(F_i\) is built exactly \(T_{i-1}\) times (once for each partial triangulation of the earlier faces). 
Therefore the total building cost across all recursive calls is
\[
    B = \sum_{i=1}^{k} T_{i-1} \cdot O(n_i),
\]
where we define \(T_0 = 1\) (the empty product). 
Using Lemma~\ref{lem:min_triangulations_final}, we have \(n_i = O(d_i)\). 
Hence
\[
    B = O\left( \sum_{i=1}^{k} T_{i-1} \cdot d_i \right).
\]
But \(T_{i-1} \cdot d_i\) is exactly the number of complete triangulations in which the triangulation of face \(F_i\) is the root triangulation (i.e., the first time we enter that face). 
Summing over all faces, we obtain \(B = O(T)\), because each complete triangulation contributes a constant number of building operations: one root triangulation per face along its unique path. 
A more formal inductive proof is given in the next section.

\section{Flipping Operations Analysis}
\label{sec:flipping_ops}

Each child triangulation in a local genealogical tree is generated by flipping a generating chord from its parent. 
The number of flips is exactly one less than the number of nodes in the tree, because the tree is connected and acyclic. 
For a given face \(F_i\), the local genealogical tree (restricted to valid triangulations compatible with \(S\)) has exactly \(d_i\) nodes (triangulations). 
Hence the number of flips performed while exploring this tree (i.e., the number of edges traversed) is \(d_i - 1\).

However, this local tree is explored multiple times: once for each triangulation of the previous faces. 
Thus the total number of flip operations for face \(F_i\) across the whole algorithm is
\[
    F_i = T_{i-1} \cdot (d_i - 1).
\]
Summing over all faces,
\[
    F = \sum_{i=1}^{k} T_{i-1} \cdot (d_i - 1) = \sum_{i=1}^{k} (T_{i-1} d_i - T_{i-1}).
\]
Now \(T_{i-1} d_i = T_i\) (the number of triangulations for the first \(i\) faces). 
Also \(\sum_{i=1}^{k} T_{i-1} \le \sum_{i=1}^{k} T_i \le k T_k = O(T)\) because each \(T_i \le T_k\) and \(k = O(n)\). 
Thus
\[
    F = \left( \sum_{i=1}^{k} T_i \right) - \left( \sum_{i=1}^{k} T_{i-1} \right) = T_k - T_0 = T - 1.
\]
Wait, this telescopes only if we consider the sum of \(T_i\) from \(i=0\) to \(k\)? Let's compute carefully:
\[
\sum_{i=1}^{k} (T_{i-1} d_i) = \sum_{i=1}^{k} T_i = T_1 + T_2 + \dots + T_k.
\]
\[
\sum_{i=1}^{k} T_{i-1} = T_0 + T_1 + \dots + T_{k-1}.
\]
Then
\[
F = (T_1 + \dots + T_k) - (T_0 + T_1 + \dots + T_{k-1}) = T_k - T_0 = T - 1.
\]
Thus the total number of flip operations is exactly \(T - 1\), independent of the face sizes! 
This is a remarkable simplification that shows the flips are perfectly amortised: each non-root triangulation is generated by exactly one flip.

\section{Base Case: Two Faces}
\label{sec:base_case_two_faces}

To illustrate the amortisation argument, consider a graph with two faces \(F_1\) and \(F_2\) having \(n_1, n_2\) vertices and generating \(d_1, d_2\) triangulations respectively. 
The total number of complete triangulations is \(T = d_1 \times d_2\).

Building operations:
\begin{itemize}
    \item Root triangulation of \(F_1\) is built once: \(O(n_1) = O(d_1)\).
    \item For each of the \(d_1\) triangulations of \(F_1\), we build a root triangulation of \(F_2\): \(d_1 \cdot O(n_2) = d_1 \cdot O(d_2) = O(T)\).
\end{itemize}
So \(B = O(T)\).

Flipping operations:
\begin{itemize}
    \item In the local tree of \(F_1\), we perform \(d_1 - 1\) flips.
    \item In the local tree of \(F_2\), for each of the \(d_1\) triangulations of \(F_1\), we perform \(d_2 - 1\) flips: total \(d_1 (d_2 - 1)\).
    \item Sum: \((d_1 - 1) + d_1(d_2 - 1) = d_1 d_2 - 1 = T - 1\).
\end{itemize}
Thus \(F = T - 1 = O(T)\).

Hence \(B + F = O(T)\).

\section{Inductive Step: From \(k-1\) to \(k\) Faces}
\label{sec:inductive_step}

We now prove by induction on the number of faces \(k\) that the total number of building and flipping operations is \(O(T)\), where \(T = \prod_{i=1}^{k} d_i\) is the total number of triangulations.

\noindent\textbf{Inductive hypothesis:} For any biconnected plane graph with at most \(k-1\) faces, the algorithm runs in \(O(T_{k-1})\) total time, where \(T_{k-1}\) is the number of triangulations of that graph.

\noindent\textbf{Inductive step:} Consider a graph with \(k\) faces \(F_1, \dots, F_k\). 
The algorithm first processes face \(F_1\) and generates all its \(d_1\) triangulations. 
For each triangulation \(t_{1,j}\) of \(F_1\), it recursively processes the remaining \(k-1\) faces (which form a subgraph with the chords of \(t_{1,j}\) added to the global set). 
By the inductive hypothesis, the time spent for a fixed \(t_{1,j}\) is \(O(T_{k-1}^{(j)})\), where \(T_{k-1}^{(j)}\) is the number of triangulations of the subgraph when \(t_{1,j}\) is fixed. 
Since the subgraphs for different \(j\) are independent (they share the same face structure but have different chord sets), the total time for the recursion is
\[
\sum_{j=1}^{d_1} O(T_{k-1}^{(j)}) = O\left( \sum_{j=1}^{d_1} T_{k-1}^{(j)} \right) = O\left( \prod_{i=2}^{k} d_i \cdot d_1 \right) = O(T_k).
\]
The building time for the root triangulations of \(F_1\) is \(O(d_1)\) (since \(n_1 = O(d_1)\)) and for the root triangulations of the remaining faces we have already accounted in the recursion. 
The flips within the local tree of \(F_1\) contribute \(d_1 - 1\) operations, which is also \(O(T)\). 
Thus the total is \(O(T)\). 
This completes the induction.

\section{Main Time Complexity Theorem}
\label{sec:time_theorem}

We now state the main result.

\begin{theorem}[Amortised constant time per triangulation]
    \label{thm:time_complexity_final}
    The algorithm generates all \(T\) triangulations of a biconnected plane graph with \(n\) vertices in \(O(T)\) total time, i.e., \(O(1)\) amortised time per triangulation.
\end{theorem}

\begin{proof}
    From the analysis above, the total number of building operations \(B\) satisfies \(B = O(T)\), and the total number of flipping operations \(F = T-1 = O(T)\). 
    Each operation takes \(O(1)\) time (hash set operations are expected \(O(1)\); list operations are worst-case \(O(1)\)). 
    Hence the total running time is \(O(T)\). 
    Since the algorithm outputs \(T\) triangulations, the amortised time per triangulation is \(O(1)\).
\end{proof}

\begin{remark}
    The hash set operations are expected constant time; if we use a deterministic balanced tree, the complexity becomes \(O(T \log n)\), which is still acceptable but not strictly \(O(T)\). 
    In practice, hash sets perform very well.
\end{remark}

\section{Space Complexity}
\label{sec:space_complexity}

We now analyse the memory usage of the algorithm.

\begin{theorem}[Linear space complexity]
    \label{thm:space_complexity_final}
    The algorithm uses \(O(n)\) space, where \(n\) is the number of vertices of the input graph.
\end{theorem}

\begin{proof}
    The algorithm stores the following data:

    \begin{enumerate}
        \item \textbf{Vertex and face information:} The input graph is stored in a doubly connected edge list (DCEL) or adjacency list, which takes \(O(n)\) space (since a planar graph has \(O(n)\) edges).
        
        \item \textbf{Global chord set \(S\):} At any point during the DFS, \(S\) contains the chords of the current partial triangulation along the recursion path. 
        The partial triangulation is a subgraph of a triangulation of a planar graph, hence it has at most \(3n-6\) edges (by Euler’s formula). 
        Therefore \(|S| = O(n)\). 
        Implementing \(S\) as a hash set uses \(O(n)\) space (each stored chord takes constant space).
        
        \item \textbf{Recursion stack:} The outer recursion depth is at most the number of faces \(k\). 
        For a biconnected plane graph, Euler’s formula gives \(k \le 2n - 4 = O(n)\). 
        The inner recursion depth (within a face’s local genealogical tree) is at most \(n_i - 2 = O(n)\). 
        However, we do not nest the inner recursion across faces; the stack frames for the inner recursion are allocated and freed for each face. 
        The maximum total stack depth is therefore \(O(n)\). 
        Each stack frame stores a constant amount of data (pointers to the current triangulation, VGS, etc.). 
        Hence the stack space is \(O(n)\).
        
        \item \textbf{Current triangulation representation:} For a face with \(n_i\) vertices, we store the lists \(L\), \(GS\), and \(VGS\), each of size \(O(n_i)\). 
        Additionally, we store opposite pairs for generating chords, also \(O(n_i)\). 
        Since we process faces one at a time, the total memory for these lists is dominated by the largest face, which has at most \(n\) vertices. 
        Thus the space for the current triangulation is \(O(n)\).
        
        \item \textbf{Backtracking information:} To restore the global set \(S\) when backtracking, we do not need to store the entire history; we can simply remove the chords that were added at each step (using the same hash set). 
        No extra space is required beyond the recursion stack.
    \end{enumerate}

    Summing all components, the total space complexity is \(O(n) + O(n) + O(n) + O(n) = O(n)\).
\end{proof}

\begin{corollary}
    The algorithm uses space linear in the number of vertices and independent of the number of triangulations \(T\), which can be exponential in \(n\).
\end{corollary}

\section{Complexity Summary}
\label{sec:complexity_summary}

Table~\ref{tab:complexity_summary_final} summarises the complexity bounds.

\begin{table}[htbp]
    \centering
    \caption{Complexity summary of the algorithm.}
    \label{tab:complexity_summary_final}
    \begin{tabular}{|l|l|}
        \hline
        \textbf{Measure} & \textbf{Complexity} \\
        \hline
        Building time per face (safe root + root triangulation + VGS) & \(O(n_i)\) \\
        Flipping time per child & \(O(1)\) \\
        Total building operations \(B\) & \(O(T)\) \\
        Total flipping operations \(F\) & \(T-1\) \\
        Total running time & \(O(T)\) \\
        \textbf{Amortised time per triangulation} & \(\mathbf{O(1)}\) \\
        Space (vertex/face data) & \(O(n)\) \\
        Space (global chord set \(S\)) & \(O(n)\) \\
        Space (recursion stack) & \(O(n)\) \\
        Space (current triangulation) & \(O(n)\) \\
        \textbf{Total space} & \(\mathbf{O(n)}\) \\
        \hline
    \end{tabular}
\end{table}

\section{Experimental Validation of Complexity}
\label{sec:complexity_experimental}

The theoretical bounds are confirmed by the experiments described in Chapter~\ref{ch:experiments}. 
For cycles of increasing size, the time per triangulation remains nearly constant (around 1.4-2.5 microseconds per triangulation), and the memory per vertex stays within a few hundred bytes. 
Large instances with tens of millions of triangulations run in time proportional to the output size, demonstrating the practical \(O(1)\) amortised behaviour.

\section{Summary}
\label{sec:complexity_summary_end}

We have proved that our algorithm runs in \(O(1)\) amortised time per generated triangulation and uses \(O(n)\) space. 
The key ingredients are the linear lower bound on the number of triangulations per face (Lemma~\ref{lem:min_triangulations_final}) and the telescoping sum that shows the total number of flips is exactly \(T-1\). 
The space analysis relies on planarity (Euler’s formula) and the fact that we never store the entire tree, only the current path. 
These complexity results match the optimal bounds achieved by the Parvez-Rahman-Nakano algorithm for triconnected graphs, while our algorithm works for the strictly larger class of biconnected graphs.